% Theme showcase. Change the theme option to compare palettes.
% Choices: classic, ocean, forest, sunset, violet.
\documentclass[theme=ocean]{CarlosTexNote}
\title{CarlosTexNote 配色主题预览}
\author{同一内容，不同语义色板}
\date{\today}
\begin{document}
\maketitle
\chapter{主题预览}
\section{数学知识块}
\begin{definition}[矩阵可逆]
若存在矩阵 $B$ 使得 $AB=BA=I$，则称方阵 $A$ 可逆，$B$ 称为 $A$ 的逆矩阵。
\end{definition}
\begin{theorem}[可逆判定]
$n$ 阶方阵 $A$ 可逆的充要条件是 $\det A\ne 0$。
\end{theorem}
\begin{property}[行列式]
$\det(AB)=\det A\det B$，且 $\det(A^{\mathsf T})=\det A$。
\end{property}
\begin{example}[行列式计算]
\[
D=\begin{vmatrix}1&2\\2&5\end{vmatrix}=1.
\]
\end{example}
\begin{solution}
按第一行展开，$D=1\cdot5-2\cdot2=1$。
\end{solution}
\section{政治知识块}
\paragraph{考点一：核心概念}
\subparagraph{基本内涵} \keyword{实践} 是人能动地改造客观世界的物质活动；理解时要抓住主体、客体和中介三个要素。
\begin{keybox}[背诵锚点]
主体是人，客体是实践对象，中介是实践工具和手段。三者共同构成实践活动的基本结构。
\end{keybox}
\begin{warningbox}[易错辨析]
不能把实践简单等同于认识活动，也不能把实践主体、客体和中介的关系割裂。
\end{warningbox}
\begin{methodbox}[复习方法]
先说清定义，再列构成要素，最后用一个具体例子解释三者关系。
\end{methodbox}
\begin{summarybox}[一句话总结]
实践是主体借助中介作用于客体的现实活动。
\end{summarybox}
\begin{recallbox}[闭卷提取]
不看笔记，能否说出实践的三个基本要素，并解释它们之间的关系？
\end{recallbox}
\end{document}
